Bestimmen Sie die Potenzreihe der Funktion
\[
f(z)
=
\begin{cases}
\displaystyle \frac{e^z - 1}{z} & \qquad\text{falls $z\ne 0$}\\
              1                 & \qquad\text{falls $z=0$}
\end{cases}
\]
an der Stelle $z=0$.
Bestimmen Sie den Konvergenzradius.

\begin{loesung}
Aus der Reihenentwicklung
\[
e^z
=
1 + z + \frac{z^2}{2!} + \frac{z^3}{3!} + \frac{z^4}{4!} + \dots
=
\sum_{k=0}^\infty \frac{z^k}{k!}
\]
der Exponentialfunktion bekommt man erst die Potenzreihe
\[
e^z - 1
=
z + \frac{z^2}{2!} + \frac{z^3}{3!} + \frac{z^4}{4!} + \dots
=
\sum_{k=1}^\infty \frac{z^k}{k!}
\]
als Potenzreihe von $e^z-1$.
Da diese Funktion offenbar eine Nullstelle bei $z=0$ hat, kann man
durch $z$ dividieren und erhält die gesuchte Potenzreihe
\[
\frac{e^z - 1}{z}
=
1 + \frac{z}{2!} + \frac{z^2}{3!} + \frac{z^3}{4!} + \dots
=
\sum_{k=1}^\infty
\frac{z^k}{(k+1)!}.
\]
Da die Funktion $f(z)$ in ganz $\mathbb{C}$ definiert ist, ist
der Konvergenzradius $\infty$.

Da alle Koeffizienten von $0$ verschieden sind, kann der
Konvergenzradius auch mit dem Quotientenkriterium gefunden werden.
Tatsächlich ist
\[
\varrho
=
\lim_{n\to\infty}
\biggl|\frac{a_n}{a_{n+1}}\biggr|
=
\lim_{n\to\infty}
\frac{1/(n+1)!}{1/(n+2)!}
=
\lim_{n\to\infty}
\frac{(n+2)!}{(n+1)!}
=
\lim_{n\to\infty}
(n+2)
=
\infty.
\qedhere
\]
\end{loesung}
